\section{P3 closure spine}
\label{sec:p3closurespine}

The detailed P3 typed-role extraction, O/Q and M/G stabilizer analysis and
cotangent-port construction are preserved in
Appendix~\ref{app:p3technical}.  The first theorem below reconstructs the
positive quotient
\[
 (E_{P3}^{\rm phys},K_*),\qquad \dim E_{P3}^{\rm phys}=7,
\]
from a smaller support--port packet.  The subsequent representation argument
then uses its canonical exterior action
\[
 \mathsf C_{\rm ext}(v)=\varepsilon(v)
 +\iota_{K_*(v,\cdot)},
\]
and the normalized two-leg record factor.  The appendix proves exactly which
complete first-jet, metric, typing and operational-port premises are required;
it does not import their physical base--seed occupation into the theorems
below.

The proof spine should be read in five steps.  First, the source-owned support
structure isolates two singleton roles and two unresolved role planes.
Second, ordered ports and the independent unit reconstruct a typed rank-seven
carrier, while the decoder preserves that rank in the body.  Third, the
class-minimal variational law occupies the reconstructed packet rather than
adding occupation as a separate postulate.  Fourth, the operator-square and
pointed-record terms produce the Clifford--record realization.  Finally, the
public pair test asks whether its source-whitened anticommutator relation can
be evaluated without terminal fitting.  Thus the sequence is
\[
 \begin{gathered}
 \text{support structure}
 \longrightarrow \text{rank-seven reconstruction}
 \longrightarrow \text{variational occupation}
 \\
 \longrightarrow \text{Clifford--record packet}
 \longrightarrow \text{empirical pair falsifier}.
 \end{gathered}
\]
The first four arrows are conditional internal results; the last is the
empirical decision frontier and currently remains data-ineligible.

\subsection{Tau-local origin of standard subsystem composition}
\label{sec:taustandardcomposition}

The two-leg record carrier used later should not be inserted merely because
standard four-dimensional notation uses a tensor product.  Let $H_0$ and
$H_1$ be finite occupied subhypergraphs of the one universal seed--body
realization.  Call them \emph{parent relation-free} when no occupied primitive
incidence meets both supports.  For this scope impose the following finite
parent-law clauses:
\begin{enumerate}[label=(LC\arabic*),nosep]
\item every primitive closure rule has connected source support contained in
  one occupied component;
\item a rule supported in $H_i$ has output support inside the component
  closure of $H_i$, so it creates neither a bridge to $H_{1-i}$ nor a shared
  external mediator;
\item the primitive component rules exhaust joint admissibility: there is no
  additional global accept/reject predicate on otherwise lawful endpoint
  pairs; and
\item the finite record linearization is source faithful and inherits the
  source-rigid unrescaled monoidal metric.
\end{enumerate}
The first three clauses concern closure-set factorization.  The fourth fixes
the common action normalization after factorization; it is not needed for the
abstract tensor shape.

\begin{theorem}[Tau-local standard subsystem composition]
\label{thm:taustandardcomposition}
Under LC1--LC3,
\begin{equation}
 \boxed{
 \mathrm{Cl}_\tau(H_0\sqcup H_1)
 \simeq
 \mathrm{Cl}_\tau(H_0)\times\mathrm{Cl}_\tau(H_1).}
 \label{eq:tauclosureproduct}
\end{equation}
Consequently the finite record algebras obey the canonical isomorphism
\begin{equation}
 \boxed{
 \mathbb C^{\mathrm{Cl}_\tau(H_0\sqcup H_1)}
 \simeq
 \mathbb C^{\mathrm{Cl}_\tau(H_0)}
 \otimes
 \mathbb C^{\mathrm{Cl}_\tau(H_1)}.}
 \label{eq:tautensorlinearization}
\end{equation}
Local parent operations have the standard commuting form
\begin{equation}
 T_0\longmapsto T_0\otimes I,\qquad
 T_1\longmapsto I\otimes T_1,\qquad
 [T_0\otimes I,I\otimes T_1]=0.
 \label{eq:taulocaloperators}
\end{equation}
Under LC4 the isomorphism in
Eq.~\eqref{eq:tautensorlinearization} is isometric in the inherited common
action unit, with no independent central tensor gain.
\end{theorem}

\begin{proof}
Take any finite global closure sequence starting from $H_0\sqcup H_1$.
LC1 assigns every primitive step to one component, and LC2 keeps its output
inside that component.  Projecting the sequence onto the two supports
therefore gives a pair of lawful local closure sequences and hence a map from
the left side of Eq.~\eqref{eq:tauclosureproduct} to the right side.

Conversely, rules on disjoint supports commute: a rule on $H_0$ changes
neither the match nor the output of a rule on $H_1$, and conversely.
Arbitrary local closure sequences can therefore be interleaved into a lawful
global sequence.  LC3 forbids a second global condition from deleting the
resulting endpoint pair.  The two constructions are inverse, proving
Eq.~\eqref{eq:tauclosureproduct}.

For finite sets $A$ and $B$, the delta-basis correspondence
\begin{equation}
 \delta_{(a,b)}\longleftrightarrow\delta_a\otimes\delta_b
 \label{eq:taudeltatensor}
\end{equation}
extends uniquely to the algebra isomorphism
$\mathbb C^{A\times B}\simeq\mathbb C^A\otimes\mathbb C^B$, proving
Eq.~\eqref{eq:tautensorlinearization}.  Component-local rules act on only one
delta factor, which gives Eq.~\eqref{eq:taulocaloperators}.  Finally LC4 makes
the primitive delta bases orthonormal in the inherited source metric and the
source-rigid monoidal continuation excludes a separate positive central
rescaling.  This proves the metric statement.
\end{proof}

\paragraph{Sharp controls and scope.}
Let both local alternative sets be $\{0,1\}$ while the global endpoint set is
only
\begin{equation}
 \mathcal D=\{(0,0),(1,1)\}.
 \label{eq:tauhiddendiagonal}
\end{equation}
Both marginal projections are onto, yet
$\dim\mathbb C^{\mathcal D}=2<4$; this hidden global relation violates LC3.
A crossing rewrite $(i,j)\mapsto(1-i,1-j)$ generates exactly the same
diagonal orbit and violates LC1--LC2.  Thus disjoint labels, complete
marginals, action valuation or a vanishing mixed Hessian do not imply
Eq.~\eqref{eq:tauclosureproduct}.  A positive central metric rescaling
preserves the tensor algebra but violates the isometric part of LC4.

The theorem is a Tau-only derivation inside the declared enriched
component-local parent class.  It does not show that the old-bare
admissibility skeleton or the unrestricted physical base--seed law realizes
LC1--LC4 on every claimed isolated scope.  If a crossing hyperedge is
occupied, the relation-free premise fails and this closure-factorization
theorem simply does not apply.  Retaining a tensor carrier as the kinematic
envelope of interacting quantum states then uses the separate strong
dagger-monoidal quantum representation result; it does not follow from
Eq.~\eqref{eq:tauclosureproduct} alone.

The full seven-role carrier need not be named in advance.  Let
$E_{\rm rel}$ be the positive six-dimensional relative source carrier and let
$P_{BMG},P_{OQ},P_{RMG}$ be commuting $K_{\rm rel}$-self-adjoint
projectors.  Write $E_{\epsilon}$ for their joint binary eigenspaces.

\begin{theorem}[Support--port reconstruction and class-minimal enriched-law realization]
\label{thm:p3supportportparent}
Assume
\begin{equation}
 \dim E_{100}=1,\qquad \dim E_{010}=2,\qquad
 \dim E_{001}=1,\qquad \dim E_{101}=2,
 \label{eq:p3supportmultiplicities}
\end{equation}
with every other joint eigenspace zero.  Assume further that the two
two-dimensional blocks possess source-owned quotient-basic ordered
cotangent decompositions
\begin{equation}
 E_{010}=L_O\oplus L_Q,
 \qquad E_{101}=L_M\oplus L_G,
 \label{eq:p3orderedports}
\end{equation}
and that the common-unit vector $z_U$ splits directly from $E_{\rm rel}$.
Then
\begin{equation}
 E_{P3}^{\rm phys}
 =L_B\oplus L_O\oplus L_R\oplus L_U
  \oplus L_Q\oplus L_M\oplus L_G,
 \qquad \dim E_{P3}^{\rm phys}=7,
 \label{eq:p3supportportrankseven}
\end{equation}
where $L_B=E_{100}$, $L_R=E_{001}$ and $L_U=\mathbb Rz_U$.
If a smooth body realization has a local left decoder, this rank is
preserved in the stabilized body.

Normalize the already stabilized body branch by
$\widehat\Gamma_{\rm body}
=\Gamma_{\rm body}-\Gamma_{\rm body}[M_*]\geq0$ on its declared stability
basin.  Write $\Phi_{\rm port}=\Gamma_{\rm port}^{\rm CTC}$ and normalize the
remaining post-body defects by
$\Phi_{P3}=\Gamma_{P3}^{\rm enr}/\mathcal A_*$ and
$\Phi_{\rm rec}=\Gamma_{\rm rec}/\mathcal A_*$.  Define the proposed
class-minimal enriched P3 law within the declared finite,
source-irredundant one-copy completion class by
\begin{equation}
 \boxed{
 \mathcal L_{P3}^{\min}
 =\left(
 \mathcal D_{\rm SF},
 \widehat\Gamma_{P3}^{\min}
 =\widehat\Gamma_{\rm body}
 +\mathcal A_*\bigl(
   \Phi_{\rm port}+\Phi_{P3}+\Phi_{\rm rec}\bigr);
 \operatorname{Real}_B,Q_B,\Omega_\varnothing;
 \mathsf{Phys}=\operatorname*{Argmin}_{\mathcal D_{\rm SF}}
 \widehat\Gamma_{P3}^{\min}
 \right),}
 \label{eq:p3minimalparentlaw}
\end{equation}
where $\mathcal D_{\rm SF}$ owns
Eq.~\eqref{eq:p3supportmultiplicities}, the unit split and the two protected
cotangent source covectors;
$\Gamma_{\rm port}^{\rm CTC}$ is the normalized two-block
Fenchel--Young law; $\Gamma_{P3}^{\rm enr}$ is the once-counted graph plus
operator-square law; $\Gamma_{\rm rec}$ is the normalized two-leg record
law; $Q_B\operatorname{Real}_B=I$; and $\Omega_\varnothing$ is the pointed
cyclic exterior root with the complete physical P3 null and no bypass.  All
three $\Phi$ terms are nonnegative dimensionless defects in the one common
action unit.  Then the physical branch selected by
Eq.~\eqref{eq:p3minimalparentlaw} satisfies
\begin{equation}
 \boxed{
 (B_\tau^{\rm phys},s_U,\mathcal L_{P3}^{\min})
 \Longrightarrow \mathfrak P_{P3}^{\min},}
 \label{eq:p3minimalparentarrow}
\end{equation}
where $\mathfrak P_{P3}^{\min}$ is the class-minimal faithful typed rank-seven
Clifford--record packet, unique up to typed source isometry and typed unitary
equivalence in the declared source-irredundant one-copy class.

The earlier narrow reduct does not entail
$\mathcal L_{P3}^{\min}$: an equal-upstream no-response completion omits a
required ordered port, while dependent-unit and equal-Gram commuting controls
independently remove the rank and operator-square conclusions.
\end{theorem}

\begin{proof}
Simultaneous orthogonal spectral decomposition of the three projectors gives
the intrinsic blocks $L_B\widehat\oplus E_{010}
\widehat\oplus L_R\widehat\oplus E_{101}$.  The two ordered cotangent
decompositions split the only repeated blocks, and the direct unit split adds
one independent line, proving Eq.~\eqref{eq:p3supportportrankseven}.
Differentiating the decoder identity gives
$DQ_BD\operatorname{Real}_B=I$, so body realization cannot collapse a typed
direction.  The two Fenchel--Young graph terms have zero response Schur
remainder.  The faithful separating operator-square term vanishes exactly on
$C(u)C(v)+C(v)C(u)=2K_*(u,v)I$, and eliminating the graph response leaves
that selector once.  The pointed exterior root and normalized record factor
then give $J_{\rm ext}^\dagger J_{\rm ext}=I_{256}$ and faithful cyclic
occupation.  The canonical exterior Clifford incidence, its graph response,
the two metric-dual ordered ports, the stabilized body and the normalized
two-leg record root form an explicit simultaneous zero of every term in
$\widehat\Gamma_{P3}^{\min}$.  Its global minimum is therefore zero.  Since
all summands are nonnegative, every global minimizer makes every summand
vanish; packet occupation follows from the declared variational law and is
not a separate common-zero premise.  Complete physical null and no bypass
give the stated
minimal-class uniqueness.  The three cited controls prove non-entailment from
the narrow reduct.  Appendix~\ref{app:p3technical} gives the componentwise
deletion certificate.  QED.
\end{proof}

Equation~\eqref{eq:p3minimalparentarrow} closes a formal ownership arrow only
after the proposed physical law is specified.  It does not prove that Nature
realizes this proposal, and it does not establish ambient-parent
exhaustivity from finite readouts.

\subsection{One common law for locality and the connected P3 packet}
\label{sec:jointlocalp3}

The locality theorem and the P3 occupation theorem act on different parts of
the source decomposition.  Let $\pi_0(s_U)$ denote the connected components
of occupied primitive incidence and define
\begin{equation}
 \boxed{
 \widehat\Gamma_{\rm CLP3}[X]
 =\sum_{C\in\pi_0(s_U)}\widehat\Gamma_C[X_C],
 \qquad K_{\rm CLP3}=\bigoplus_C K_C.}
 \label{eq:jointlocalp3action}
\end{equation}
Assume that primitive rules exhaust their full connected supports, preserve
component closure, and admit no additional cross-component multiplier or
hidden endpoint predicate; that every component action is nonnegative,
once-counted and has an attainable zero; and that realization is source
faithful in the common unrescaled action unit.  On the occupied connected
component $C_{P3}$ take $\widehat\Gamma_C$ to be the already defined
$\widehat\Gamma_{P3}^{\min}$ of
Eq.~\eqref{eq:p3minimalparentlaw}.

\begin{theorem}[Component-local P3 joint occupation]
\label{thm:jointlocalp3}
Under these hypotheses every global minimizer of
Eq.~\eqref{eq:jointlocalp3action} satisfies LC1--LC4 on every
parent-relation-free support pair and occupies the complete typed rank-seven
packet $\mathfrak P_{P3}^{\min}$ on $C_{P3}$.  The earlier narrow
base--seed/body reduct does not entail this common law.
\end{theorem}

\begin{proof}
Full-support assignment and component preservation give LC1--LC2.  Primitive
exhaustivity and absence of an extra cross-component multiplier give LC3,
while source-faithful unrescaled realization is LC4.  Hence
Theorem~\ref{thm:taustandardcomposition} applies to every relation-free pair.
Nonnegativity, disjoint component variables and attainability give
\begin{equation}
 \operatorname*{Argmin}\widehat\Gamma_{\rm CLP3}
 =\prod_C\operatorname*{Argmin}\widehat\Gamma_C,
 \qquad \min\widehat\Gamma_{\rm CLP3}=0.
 \label{eq:jointlocalp3argmin}
\end{equation}
Restriction to $C_{P3}$ is therefore a zero of
$\widehat\Gamma_{P3}^{\min}$, so
Theorem~\ref{thm:p3supportportparent} forces the complete rank-seven
Clifford--record packet.  Its internal incidences do not violate LC1--LC2:
they belong to one connected component and are never factorized by the
theorem.

For non-entailment, keep the local rules and complete marginals but impose the
hidden diagonal endpoint relation in Eq.~\eqref{eq:tauhiddendiagonal}; LC3
fails.  Independently replace the P3 realization by the equal-total-scalar
rank-six missing-port control used in the support--port deletion certificate.
Their product has the same earlier narrow descriptors while failing both
targets.  Thus that reduct cannot select
Eq.~\eqref{eq:jointlocalp3action}.
\end{proof}

This theorem closes compatibility and common-law occupation inside the
declared enriched completion.  It is not a proof that Nature realizes
$\widehat\Gamma_{\rm CLP3}$ or that the listed components exhaust the ambient
parent.

\subsection{Why the class-minimal law is not a terminal fit}

Here \emph{class-minimal} means componentwise minimal only within the
declared finite, source-irredundant one-copy completion class.  In both the
law and packet symbols, the superscript ``min'' has this restricted meaning;
it does not assert
that the law is Nature's unique or globally minimal parent law.  Nor is the
functional fitted backwards from terminal residuals: its inputs are frozen
on the source side before endpoint evaluation.  The origin of its clauses is
as follows.
\begin{enumerate}[label=(\roman*),nosep]
\item The support multiplicities and projectors are inherited from the
  source-role structure.  Their realization by the unrestricted physical
  base--seed system remains an explicit completion premise.
\item The O/Q and M/G cotangent ports are minimal preimages, within this
  class, of already required ordered dualities.  Deleting either port leaves
  an explicit typed swap symmetry.
\item The independent common unit and the left decoder have separate rank
  roles: the unit raises source rank from six to seven, while the decoder
  prevents collapse in body realization.  Neither is a terminal-specific
  gain.
\item The state-neutral operator-square term is the genuinely new
  Tau-specific cross-role clause.  Without it, a commuting packet can retain
  the same normalized marginal trace Gram.
\item The pointed cyclic exterior root and no-bypass clause are required by
  the faithful full-rank nontracial and extra-record countermodels.  They
  select record occupation and completeness only within the declared class.
\item The two-leg record tensor is not inferred from a terminal residual.
  Theorem~\ref{thm:taustandardcomposition} derives it from the independent
  component-local no-bridge closure law; the diagonal hidden-relation model
  prevents replacing that source law by complete marginals or action
  additivity.
\end{enumerate}
Thus Theorem~\ref{thm:p3supportportparent} is a noncircular internal
realization theorem for a source-frozen proposed law.  It is not a deeper
derivation of that law from the narrow parent, and it does not exclude richer
parent laws with the same physical zero locus.

\begin{theorem}[Pointed exterior-seed closure and rank-only no-go]
\label{thm:p3exteriorseedclosure}
Assume that the occupied typed P3 seed owns the complete oriented word
grammar $\Lambda^\bullet E_{P3}$, explicitly occupies the empty word
$\Omega_\varnothing$, and that the base realizes
$\mathsf C_{\rm ext}(v)$ as in Eq.~\eqref{eq:p3extclifford}.  Assume also that
the normalized two-leg Frobenius record root generates an orthonormal record
basis $r_0,r_1$ with no independent bypass.  Then the vectors
\begin{equation}
 \Psi_{I,a}
 =\mathsf C_{\rm ext}(e_{i_k})\cdots
  \mathsf C_{\rm ext}(e_{i_1})\Omega_\varnothing\otimes r_a,
 \qquad I=\{i_1<\cdots<i_k\}\subseteq\{1,\ldots,7\},
 \quad a\in\{0,1\},
 \label{eq:p3pointedorbit}
\end{equation}
form an orthonormal basis of $\mathcal F_{P3}$.  Equivalently, the root-orbit
map $J_{\rm ext}:\mathbb C^{256}\to\mathcal F_{P3}$ obeys
\begin{equation}
 \boxed{J_{\rm ext}^{\dagger}J_{\rm ext}=I_{256}.}
 \label{eq:p3exteriorgramlock}
\end{equation}
Consequently the relation, source-symmetry and cyclic trace--GNS conditions
are simultaneous kinematic consequences on this declared carrier; no
independent Gram-isometry action term is required there.

Full orbit rank, source nonloss and equal central block weights are not
sufficient substitutes.  On
$\mathcal A=M_8(\mathbb C)^{\oplus4}$ with normalized trace, choose
\begin{equation}
 h_\epsilon
 =\bigoplus_{b=1}^{4}
  \operatorname{diag}(1+\epsilon,1-\epsilon,1,\ldots,1),
 \qquad 0<\epsilon<1,
 \label{eq:p3nontracialdensity}
\end{equation}
and $J_\epsilon(a)=a h_\epsilon^{1/2}$ on the regular carrier.  Then
$J_\epsilon$ has rank 256 and equal weight $1/4$ on every central block, but
\begin{equation}
 J_\epsilon^\dagger J_\epsilon=R_{h_\epsilon}\ne I,
 \qquad
 \omega_\epsilon([E_{01}^{(1)},E_{10}^{(1)}])=\frac{\epsilon}{16}\ne0.
 \label{eq:p3rankonlycontrol}
\end{equation}
Thus this faithful full-rank root state is nontracial and retains a nonzero
metric defect.
\end{theorem}

\begin{proof}
Creation followed by $K_*$-adjoint contraction sends the empty exterior word
to the signed wedge basis vector labelled by the same subset $I$.  Distinct
subsets are orthogonal, and every subset occurs once.  Tensoring with the two
orthonormal record vectors therefore gives exactly $2^7\cdot2=256$
orthonormal vectors, proving Eq.~\eqref{eq:p3exteriorgramlock}.  The Clifford
relations give relation generation, the induced exterior action transports
every source isometry, and the full word orbit is cyclic; the normalized
regular vector state is the represented trace.

For the control, positivity of $h_\epsilon$ makes right multiplication by
$h_\epsilon^{1/2}$ invertible, so the orbit map has full rank.  Every block
has trace eight, hence the four central weights are equal.  Direct evaluation
in the normalized regular inner product gives
$J_\epsilon^\dagger J_\epsilon=R_{h_\epsilon}$.  Finally
$[E_{01},E_{10}]=E_{00}-E_{11}$, whose $h_\epsilon$-weighted normalized trace
in one of the four blocks is $2\epsilon/32=\epsilon/16$.  The root state is
therefore not tracial and its Gram cannot be the source identity. QED.
\end{proof}

The positive half of Theorem~\ref{thm:p3exteriorseedclosure} is conditional
on physical ownership of the complete pointed exterior grammar.  The
negative half is an exact countermodel: dimension, rank and the retained
discrete block symmetries cannot replace that source-owned pointed metric.
The unrestricted physical base--seed law has not yet been shown to own this
grammar.


Appendix~\ref{app:p3certificate} packages the same requirement as an exact
complete-incidence characterization: within the declared finite
source-irredundant Clifford--recovery class, a minimal faithful typed P3
subpacket exists uniquely up to typed unitary equivalence if and only if the
parent owns the complete rank-seven incidence.  This is a characterization of
the missing object, not a proof of its unrestricted physical existence.


The certificate also admits one finite positive source selector.  On the
dimensionless normalized quotient $(E_{P3}^{\rm phys},K_*)$, let
$n=\dim E_{P3}^{\rm phys}=7$, let $C:E_{P3}^{\rm phys}\to
\mathcal A_{\rm body}^{\rm sa}$ be linear, and let $\tau$ be a faithful
normalized trace.  Define
\begin{equation}
 \mathcal V_{\rm OS}[C]
 = \frac{\mathcal A_*}{2}
 \int_{S(E_{P3},K_*)}
 \tau\!\left[
 (C(v)^2-I)^\dagger(C(v)^2-I)
 \right]d\mu(v),
 \label{eq:p3operatorsquarefunctional}
\end{equation}
where $d\mu$ is normalized orthogonal-invariant sphere measure.  Before
unit-sphere restriction the defect is $C(v)^2-K_*(v,v)I$, so the common
action unit $\mathcal A_*$ restores dimensions without a role-dependent gain.
For a $K_*$-orthonormal frame, write $A_a=C(e_a)$,
$d_a=A_a^2-I$, and $h_{ab}=A_aA_b+A_bA_a$.  The fourth sphere moments give
the exact finite formula
\begin{equation}
 \frac{2\mathcal V_{\rm OS}[C]}{\mathcal A_*}
 =\frac{1}{n(n+2)}\left[
 2\sum_a\lVert d_a\rVert_\tau^2
 +\left\lVert\sum_a d_a\right\rVert_\tau^2
 +\sum_{a<b}\lVert h_{ab}\rVert_\tau^2
 \right].
 \label{eq:p3operatorsquarefinite}
\end{equation}

\begin{theorem}[Operator-square source selector]
\label{thm:p3operatorsquareselector}
For a faithful normalized trace,
$\mathcal V_{\rm OS}[C]\geq0$, with equality if and only if
\begin{equation}
 C(v)^2=K_*(v,v)I
 \quad\hbox{for every }v\in E_{P3}^{\rm phys}.
 \label{eq:p3operatorsquarezero}
\end{equation}
Equivalently,
$C(u)C(v)+C(v)C(u)=2K_*(u,v)I$.  The zero set is invariant under every
$K_*$-orthogonal source-frame change.
\end{theorem}

\begin{proof}
The integrand in Eq.~\eqref{eq:p3operatorsquarefunctional} is continuous and
nonnegative.  Faithfulness of $\tau$ makes a zero integral equivalent to a
zero operator defect on the unit sphere; homogeneity extends the result to
all $v$.  Polarization gives the anticommutator, and the reverse implication
is immediate.  Orthogonal invariance follows from invariance of $K_*$ and
$d\mu$.  Finally,
$\mathbb E[v_a^4]=3/[n(n+2)]$,
$\mathbb E[v_a^2v_b^2]=1/[n(n+2)]$, and the odd moments vanish, yielding
Eq.~\eqref{eq:p3operatorsquarefinite}.
\end{proof}

This selector contains information absent from the scalar Gram matrix.  The
canonical Clifford family and seven commuting four-qubit Pauli-$Z$ words act
on the same 16-dimensional carrier and both obey
$\tau(A_aA_b)=\delta_{ab}$.  The finite certificate nevertheless gives
\begin{equation}
 \frac{2\mathcal V_{\rm OS}^{\rm Clifford}}{\mathcal A_*}=0,
 \qquad
 \frac{2\mathcal V_{\rm OS}^{\rm commuting}}{\mathcal A_*}=\frac{4}{3}.
 \label{eq:p3operatorsquarecontrol}
\end{equation}
Thus scalar Hessian/Gram data cannot select the physical packet, whereas the
operator-valued source term can.

Appendix~\ref{app:p3ownership} gives the staged enriched-parent functional
whose nonnegative zero locus combines this operator-square law with the P3
handoff without double counting.  That theorem proves ownership inside the
declared enriched completion.  It does not prove that the unrestricted
physical base--seed law realizes or occupies the enrichment.

\subsection{Handoff to the joint-terminal assembly}

The results above close only the algebra/representation level. Foundation
Paper VII-A separately owns the once-counted terminal functional, the Born
effect derivative and the conditional Einstein--Maxwell--matter equations.
Combining VII-A with Theorem~\ref{thm:p3supportportparent} is therefore a
cross-paper dependency composition, not another theorem proved in this
technical paper. In particular, rank-seven typing or nonzero activation must
not be mistaken for unrestricted physical occupation of the complete packet.


\section{Class-minimal enriched-law closure and empirical frontier}
\label{sec:selectionfrontier}

The earlier narrow reduct has the exact non-entailment
\begin{equation}
 \boxed{
 (B_\tau^{\rm phys},s_U,\mathcal L_{\rm narrow})
 \not\Longrightarrow
 \mathfrak P_{P3}^{\min}.
 }
 \label{eq:narrowparentnogo}
\end{equation}
The no-response port control already proves Eq.~\eqref{eq:narrowparentnogo};
the dependent-unit, equal-Gram commuting and nontracial-root controls locate
three further independent failures.  It is therefore impossible to remove
this no-go by another manipulation of the unchanged reduct.

Theorem~\ref{thm:p3supportportparent} instead supplies an explicit proposed
physical law and proves
\begin{equation}
 \boxed{
 (B_\tau^{\rm phys},s_U,\mathcal L_{P3}^{\min})
 \Longrightarrow \mathfrak P_{P3}^{\min}.
 }
 \label{eq:decisiveparentarrow}
\end{equation}
Thus the formal parent-law realization and occupation arrow is closed inside
the declared enriched physical MVP.  The phrase ``Nature realizes this law''
is not an additional conclusion that can be proved from an unspecified
parent: it is the empirical status of the proposed law.  Likewise,
Eq.~\eqref{eq:decisiveparentarrow} does not establish that the selected
connected packet exhausts every ambient source extension.

Two scientific routes remain, but they have different status:
\begin{enumerate}
\item \emph{Deeper-origin route.}  Derive
  $\mathcal L_{P3}^{\min}$ from a more primitive enriched base--seed law.  This
  could strengthen the theory, but is no longer required for the internal
  consistency of the stated completion; Eq.~\eqref{eq:narrowparentnogo}
  forbids obtaining it from the unchanged narrow reduct.  The standard
  subsystem-composition part of this route is now closed conditionally by
  Theorem~\ref{thm:taustandardcomposition}; the remaining deeper-origin debt
  concerns physical realization of LC1--LC4 and of the full support--port,
  operator-square and pointed-root law.
\item \emph{Discriminator route.}  Freeze a source-side morphological
  descriptor and derive a gain-free, parameter-free relation between at least
  two terminal readouts that is not implied by the corresponding standard
  comparator.  The immediate algebraic target is
  $\{C_a,C_b\}/2=\delta_{ab}I$ in a source-whitened frame; the stronger
  $\mathcal O_{P1|P3}=0$ requires an independently certified global
  source-completeness premise.  Any endpoint prediction must be fixed before
  data are scored.
\end{enumerate}
Failure of a source-frozen discriminator demotes the proposed physical law.
Neither route may be repaired by adding another standard-limit recovery or by
selecting source structure from terminal residuals.
